% GENERATED FILE. Edit canonical lessons, then run npm run generate:books.
% canonical-source-hash: fnv1a64:c3861aaffb584f30
% canonical-lessons: ES-C340-cerca, ES-C340-lejos, ES-C340-derecha, ES-C340-izquierda, ES-C340-enfrente

\chapter{Near, Far, and Which Side}
\label{ch:cerca-y-lejos}
\hlchaptermodality{340}

\begin{chapteropening}
\textbf{By the end of this chapter:} \emph{I can say cerca, lejos, a la derecha, a la izquierda and enfrente, and say why the left-hand side is the one Spanish word in the set that is older than Rome.}

Complete ES-C340-enfrente: place a thing near or far, send somebody right or left, and put one thing facing another.
\end{chapteropening}

\section[cerca]{cerca — near, from the ring drawn around you}

\label{lesson:ES-C340-cerca}



\begin{quote}
Say \emph{heat}, then \emph{age}. (\emph{El calor}; \emph{la edad}.)
\end{quote}

\subsection*{What to know first}
\begin{itemize}
\raggedright
  \item ¿dónde está? — the question this answers.
  \item ahí · allí — the two distances you already have.
\end{itemize}

\begin{sounds}
\begin{itemize}
\raggedright
  \item \emph{SER-ka}, two beats with the weight on the first. The \emph{c} before \emph{e} is soft — an \emph{s} across the Americas. $\to$ pronunciation reference
\end{itemize}
\end{sounds}

\begin{cousinweb}
\textbf{cerca} — \textquotedblleft{}near,\textquotedblright{} \textquotedblleft{}nearby.\textquotedblright{}

Latin \textbf{circā} meant \emph{around, round about} — the preposition of a ring drawn about a thing. What lies around you lies close, and Spanish let the word slide from \emph{around} to \emph{near}.

English took \textbf{circa} whole and left it in Latin, where it still marks a date given loosely: \emph{circa 1500}, somewhere in the ring around that year. The same ring is in \textbf{circle}, in a \textbf{circuit} that goes round and comes back, and in the \textbf{circus}, which was a ring long before it was a show.

So when you say \emph{cerca}, the picture underneath is a circle with you at the centre, and everything inside it counted as near.
\end{cousinweb}

\begin{grammarlens}[title={What you've built}]
The first of a handful of words for where a thing is.
\end{grammarlens}

\subsection*{Your turn}
Say these aloud:

\begin{itemize}
\raggedright
  \item \emph{cerca}
  \item \emph{muy cerca}
  \item \emph{la plaza está cerca}
\end{itemize}

\begin{tcolorbox}[breakable,colback=teal!4,colframe=teal!35!black,title={Before you move on}]
What did \emph{circā} mean? (Around, round about.) And what does English \emph{circa} still do to a date? (Marks it as roughly that year.)
\end{tcolorbox}
\section[lejos]{lejos — far, from a rope let out further}

\label{lesson:ES-C340-lejos}



\begin{quote}
Say \emph{age}, then \emph{near}. (\emph{La edad}; \emph{cerca}.)
\end{quote}

\subsection*{What to know first}
\begin{itemize}
\raggedright
  \item cerca — the other end of this pair.
\end{itemize}

\begin{sounds}
\begin{itemize}
\raggedright
  \item \emph{LE-hos}, the weight on the first beat. The \emph{j} is the throaty \emph{jota}, further back than an English \emph{h}. $\to$ pronunciation reference
\end{itemize}
\end{sounds}

\begin{cousinweb}
\textbf{lejos} — \textquotedblleft{}far,\textquotedblright{} \textquotedblleft{}far away.\textquotedblright{}

Latin \textbf{laxus} meant \emph{loose, slack, not drawn tight} — said of a rope, a rein, a bowstring. Its comparative \textbf{laxius} meant \emph{more loosely, further apart}, and Spanish froze that comparative into a plain adverb of distance.

The slack rope is all over English. Something \textbf{lax} is loose in its discipline. To \textbf{relax} is to let out again. To \textbf{release} is to slacken a hold entirely. Even a \textbf{leash} is the loose length you let a dog run on.

Two words, two pictures: \emph{cerca} draws a ring and stays inside it; \emph{lejos} pays the rope out until you are past the far end.
\end{cousinweb}

\begin{grammarlens}[title={What you've built}]
A pair now: \emph{cerca} and \emph{lejos}, the two answers to how far.
\end{grammarlens}

\subsection*{Your turn}
Say these aloud:

\begin{itemize}
\raggedright
  \item \emph{lejos}
  \item \emph{muy lejos}
  \item \emph{la casa está lejos}
\end{itemize}

\begin{tcolorbox}[breakable,colback=teal!4,colframe=teal!35!black,title={Before you move on}]
What did \emph{laxus} mean? (Loose, slack.) And which English verb means to let that slack out again? (\emph{Relax}.)
\end{tcolorbox}
\section[la derecha]{la derecha — the right hand, from the line that goes straight}

\label{lesson:ES-C340-derecha}



\begin{quote}
Say \emph{near}, then \emph{far}. (\emph{Cerca}; \emph{lejos}.)
\end{quote}

\subsection*{What to know first}
\begin{itemize}
\raggedright
  \item el camino — what you will be sending somebody along.
  \item la esquina — where the turn happens.
\end{itemize}

\begin{sounds}
\begin{itemize}
\raggedright
  \item \emph{de-RE-cha}, three beats with the weight on the middle. The \emph{ch} is one sound, as in English \emph{church}. $\to$ pronunciation reference
\end{itemize}
\end{sounds}

\begin{cousinweb}
\textbf{la derecha} — \textquotedblleft{}the right,\textquotedblright{} the right-hand side.

Latin \textbf{dīrigere} meant \emph{to set straight, to line up} — \emph{dis-} apart plus \textbf{regere}, to rule or guide. Its past participle \textbf{dīrectus} meant \emph{straightened}, and that is the word underneath \emph{derecha}.

Most languages name the right hand for what it does well, and Latin thought of it as the hand that goes straight to the mark. The left got the leftovers.

The straightening is loud in English. Something \textbf{direct} goes the straight way. A \textbf{direction} is the straight line you are pointed along. Something \textbf{erect} has been straightened upright. And to \textbf{address} a letter is to set it straight toward somebody.

Spanish keeps both jobs: \textbf{a la derecha} is \emph{to the right}, and \textbf{derecho} on its own is \emph{straight ahead} — the older meaning, still working.
\end{cousinweb}

\begin{grammarlens}[title={What you've built}]
Now you can send somebody one way at the corner.
\end{grammarlens}

\subsection*{Your turn}
Say these aloud:

\begin{itemize}
\raggedright
  \item \emph{la derecha}
  \item \emph{a la derecha}
  \item \emph{la calle está a la derecha}
\end{itemize}

\begin{tcolorbox}[breakable,colback=teal!4,colframe=teal!35!black,title={Before you move on}]
What did \emph{dīrectus} mean? (Straightened, set straight.) And what does \emph{derecho} mean when it stands alone? (Straight ahead.)
\end{tcolorbox}
\section[la izquierda]{la izquierda — the left hand, the word that is older than Rome}

\label{lesson:ES-C340-izquierda}



\begin{quote}
Say \emph{far}, then \emph{the right}. (\emph{Lejos}; \emph{la derecha}.)
\end{quote}

\subsection*{What to know first}
\begin{itemize}
\raggedright
  \item la derecha — the side this one answers.
\end{itemize}

\begin{sounds}
\begin{itemize}
\raggedright
  \item \emph{is-KYER-da}, the weight on the second beat. The \emph{z} is an \emph{s} across the Americas, and \emph{-quie-} is one gliding syllable. $\to$ pronunciation reference
\end{itemize}
\end{sounds}

\begin{cousinweb}
\textbf{la izquierda} — \textquotedblleft{}the left,\textquotedblright{} the left-hand side.

Nearly every Spanish word you have met came from Latin. This one did not.

\textbf{Izquierda} is \textbf{ezkerra}, from Basque — a language already old in the peninsula when the Romans arrived, related to nothing else alive. A few Basque words survived the whole Roman centuries and are still in daily Spanish, and this is the best known of them.

What it pushed aside was Latin \textbf{sinistra}, the ordinary Latin word for \emph{left}. Latin had long treated the left as the unlucky side, and \emph{sinistra} carried so much of that dread that Spanish quietly replaced it with a neutral borrowed word.

English went the other way and kept the discarded one. \textbf{Sinister} in English means threatening or evil, which is what Latin had already made of \emph{left} — so Spanish and English ended up with opposite halves of the same old quarrel: Spanish took the plain Basque word, English took the dark Latin one.

\textbf{A la izquierda} is \emph{to the left}.
\end{cousinweb}

\begin{grammarlens}[title={What you've built}]
Both hands now, and one word in your Spanish older than Rome.
\end{grammarlens}

\subsection*{Your turn}
Say these aloud:

\begin{itemize}
\raggedright
  \item \emph{la izquierda}
  \item \emph{a la izquierda}
  \item \emph{a la derecha o a la izquierda}
\end{itemize}

\begin{tcolorbox}[breakable,colback=teal!4,colframe=teal!35!black,title={Before you move on}]
Which language did \emph{izquierda} come from? (Basque.) And which Latin word did it push out, that English kept? (\emph{Sinistra} — English \emph{sinister}.)
\end{tcolorbox}
\section[enfrente]{enfrente — opposite, standing forehead to forehead}

\label{lesson:ES-C340-enfrente}



\begin{quote}
Say \emph{the right}, then \emph{the left}. (\emph{La derecha}; \emph{la izquierda}.)
\end{quote}

\subsection*{What to know first}
\begin{itemize}
\raggedright
  \item cerca — one of the distances this refines.
  \item la calle — what you will be crossing.
\end{itemize}

\begin{sounds}
\begin{itemize}
\raggedright
  \item \emph{en-FREN-te}, the weight on the middle beat. Both \emph{e} sounds are short and even, never sliding the way an English \emph{e} does. $\to$ pronunciation reference
\end{itemize}
\end{sounds}

\begin{cousinweb}
\textbf{enfrente} — \textquotedblleft{}opposite,\textquotedblright{} \textquotedblleft{}across from,\textquotedblright{} \textquotedblleft{}facing.\textquotedblright{}

It is two pieces you can see: \textbf{en} plus \textbf{frente}.

\textbf{Frente} is Latin \textbf{frontem}, from \textbf{frōns} — and the Latin word meant \emph{the forehead}. The forehead is the part of you that goes first, the flat face you present to whatever is coming, so \emph{frons} widened to mean the front of anything at all.

English is full of that brow. The \textbf{front} of a building is its forehead. A \textbf{frontier} is the forehead of a country, the part that faces outward. To \textbf{confront} somebody is to bring foreheads together. An \textbf{affront} is a blow struck at the forehead, which is why it means an insult. And \textbf{effrontery} is behaving as though you had no forehead left to blush with.

So \emph{enfrente} is not merely \emph{near}. It is the thing whose forehead is turned toward yours — across the street, facing you.
\end{cousinweb}

\begin{grammarlens}[title={What you've built}]
Five words for where a thing is: \emph{cerca}, \emph{lejos}, \emph{la derecha}, \emph{la izquierda}, \emph{enfrente}. Enough to place almost anything.
\end{grammarlens}

\subsection*{Your turn}
Say these aloud:

\begin{itemize}
\raggedright
  \item \emph{enfrente}
  \item \emph{enfrente de la casa}
  \item \emph{está enfrente, muy cerca}
\end{itemize}

\begin{tcolorbox}[breakable,colback=teal!4,colframe=teal!35!black,title={Before you move on}]
What did \emph{frontem} mean in Latin? (The forehead.) And what is an \emph{affront}, word for word? (A blow at the forehead.)
\end{tcolorbox}
